%%%%

%%%   Самарский государственный технический университет

%%%   Кафедра прикладной математики и информатики

%%%

%%%   Преамбула для статьи в журнал "Вестник Самарского государственного

%%%   технического университета. Серия: «Физико-математические науки»"

%%%   Ориентирована под MikTEx 2.4 for Win32

%%%

%%%   Хотя сам журнал собирается под GNU/Linux на дистрибутиве TeXLive



\documentclass[11pt,a4paper,twoside]{article}



\usepackage[cp1251]{inputenc}

\usepackage[T2A]{fontenc}

\usepackage[english,russian]{babel}

\usepackage{samgtubib}

\usepackage[dvips]{graphicx}

\usepackage[multiple]{footmisc}



\usepackage{samgtu}

\renewcommand{\VestnikYear}{2009} % Год издания Вестника

\renewcommand{\VestnikNumber}{2\,(19)} % Номер Вестника

\renewcommand{\VestnikMonth}{Ноябрь} % Месяц выхода Вестника

\renewcommand{\VestnikData}{01.11.09} % Дата подписания Вестника в печать

\renewcommand{\VestnikUPL}{26,64} % Усл. печ. л.

\renewcommand{\VestnikUIL}{25,73} % Уч.-изд. л.

\renewcommand{\VestnikRegNumber}{479/09} % Регистрационный номер

\renewcommand{\VestnikOrder}{994} % Номер заказа







%

%\begin{document}



\renewcommand{\firstpage}{119}

\renewcommand{\lastpage}{128}



\udk{544.032.65:577.1} \msc{78A70}



\title{Медико"=биологические аспекты лазерного воздействия}

{Медико-биологические аспекты лазерного воздействия}



\author{Д.~М.~Гуреев}{Д.~М.~Г\,у\,р\,е\,е\,в}



\institute{\samgtu.}



\email{E-mail}{anton\_gureev@samaradom.ru}



\otherinf{\fullauthor{Дмитрий Михайлович Гуреев}\regalia{д.ф.-м.н., проф.}\position{профессор} \dept{каф. общей физики и физики нефтегазового производства} }



\keywords{лазерное излучение, коэффициенты отражения и пропускания, биологические ткани}



\workabstract{Экспериментально изучен характер распределения коэффициентов отражения и пропускания излучения различных длин волн биологическими тканями. 

Дано физическое обоснование подходов к оптимизации выбора лазерных источников и условий их практического применения для решения фундаментальных и прикладных задач медико"=биологических исследований.}



\engtitle{Medical-and-biological aspects of laser influence}



\engauthor{D.~M.~Gureev}{D.~M.~G\,u\,r\,e\,e\,v}





\enginstitute{\engsamgtu.}



\otherenginfm{\fullauthor{Dmitriy M. Gureev} \regalia{Dr. Phys. \& Math. Sci.} \position{Professor} 

\dept{Dept. of General Physics and Physics of Oil and Gas Production}}



\engabstract{The distributions of reflection and admission coefficients of a different wave-length radiation by the biological tissues were experimentally studied. A physical basis of the methods of approach to an optimum selection of the laser sources and the conditions of its practical application to solution of the fundamental and applied problems of the medical-and-biological investigations was given.}



\engkeywords{laser radiation, reflection and admission coefficients, biological tissues}



\date{04/XII/2012}

\revisiondate{25/II/2013}





\maketitle



\Section[n]{Введение} В настоящее время широкое развитие получила

лазерная медицина, в основе которой лежит эффективное

использование уникальных свойств лазерного излучения, таких как

высокие монохроматичность, направленность, яркость, в хирургии,

терапии, диагностике \cite{gur:bib:1, gur:bib:2}. В

медико"=биологических исследованиях важное место также отводится

изучению влияния лазерного излучения на наследственность растений

с целью повышения урожайности сельскохозяйственных культур

\cite{gur:bib:3}. В лазерной терапии, диагностике, селекции

растений преимущественно используется излучение гелий"=неонового

лазера (длина волны излучения 0,633 мкм "--- красная область

спектра) низкой интенсивности $\sim 10^{-1}$~Вт/см$^2$,

сопоставимой с интенсивностью излучения Солнца во всей его широкой

области спектра на поверхности Земли в ясный солнечный день.

Вместе с тем имеются также примеры успешного практического

применения излучения аргонового лазера (длина волны излучения

0,488 и 0,515~мкм "--- сине"=зеленая область спектра), в частности

для лечения сосудистой патологии. Основу последнего составляет

избирательный характер поглощения гемоглобином крови излучения

сине"=зеленой области спектра. Выбор гелий"=неонового лазера в

качестве основного инструмента в медико"=биологических

исследованиях с использованием лазерного излучения обосновывается

наибольшей глубиной проникновения излучения красной и ближней

инфракрасной областей спектра в биологические ткани. Однако при

этом несколько в стороне остается вопрос: насколько оправданным

является такой выбор с точки зрения эффективности взаимодействия

излучения данных областей спектра с биологическими тканями?



Целью данной работы является экспериментальное изучение характера

распределения коэффициентов отражения и пропускания излучения

различных длин волн биологическими тканями и физическое

обоснование подходов к оптимизации выбора лазерных источников и

условий их практического применения в медико"=биологических

исследованиях.



\newpage



\Section[n]{Материалы и методики эксперимента} В качестве

исследуемых биологических тканей использовались срезы большинства

плодовых и овощных культур, разнообразные ягодные и зерновые

культуры, кисть руки человека, а также разновидности почв

(чернозем, песок, глина и их смеси), вода. Коэффициенты отражения

измерялись на фотометре отражения ФО-1 в диапазоне длин волн

излучения 0,364--0,927~мкм. Оптическая схема фотометра отражения

приведена на рис.~\ref{gur:pict1}. Принцип её работы состоит в следующем.

Излучение от источника света 1 (галогенной лампы КГМ 6,3-15)

попадает на конденсор 2, который переносит изображение источника в

плоскость полевой диафрагмы 3. Полевая диафрагма 3 устраняет

влияние бликов колбы источника света. После полевой диафрагмы 3

световой пучок проходит сменную диафрагму 4, изображение которой с

помощью объектива 5 и зеркала 10 переносится в плоскость

измеряемого образца 11. В зависимости от размера измеряемого

образца диаметр светового пятна в его плоскости может изменяться и

принимать значения 30, 18 и 12~мм.



\begin{figure}[b!] \centering





\includegraphics[scale=.85]{gur1}



\caption{Оптическая схема фотометра отражения ФО-1: \scriptsize 1 "---

источник света; 2 "--- конденсор; 3 "--- полевая диафрагма; 4 "---

сменные диафрагмы; 5 "--- объектив; 6 "--- светофильтры; 7, 11

"--- измеряемые образцы; 8 "--- фотометрический шар; 9 "--- экран;

10, 13 "--- зеркала; 12 "--- светоловушка; 14, 17 "--- молочные

стёкла; 15, 16 "--- приёмники излучения \label{gur:pict1}}

\end{figure}



Зеркало $10$, расположенное внутри фотометрического шара $8$,

может поворачиваться вокруг оси и занимать три фиксированных

положения: 

\begin{itemize}

\item[1)] ка\-либ\-ров\-ка-1 "--- световой поток отражается от

плоской поверхности зеркала на внутреннюю стенку шара; 

\item[2)]

калибровка-2 "--- световой поток отражается от сферической

поверхности зеркала на внутреннюю стенку шара; 

\item[3)] измерение "---

световой поток отражается плоской поверхностью зеркала в плоскость

измеряемого образца 11 или светоловушку 12.

\end{itemize}



 Свет, рассеянный в

шаре, поступает на приемники излучения 15, 16 через окна, закрытые

молочными стеклами 14, 17. В качестве приемников излучения служат

фотоумножитель ФЭУ-4 и фотодиод ФД-24К для областей спектра

0,364--0,750~мкм и 0,750--0,927~мкм соответственно.



В случае измерения коэффициента пропускания диффузно рассеивающих

образцов 7 вводится экран 9, который предохраняет приемники

излучения от попадания на них рассеянного непосредственно

образцами света.



Для выделения узких участков спектра поочередно вводятся

светофильтры 6. В состав фотометра отражения входят двенадцать

светофильтров, из них один для ближней ультрафиолетовой области

спектра (0,364~мкм), шесть для видимой области спектра (0,400;

0,457; 0,490; 0,520; 0,582; 0,620~мкм), четыре для ближней

инфракрасной области спектра (0,750; 0,832; 0,874;\linebreak

0,927~мкм) и один корригирующий светофильтр для приведения

спектральной чувствительности фотоумножителя к относительной

видности глаза.



Зеркало 13 установлено для наблюдения за положением измеряемого

образца.



Для измерения коэффициентов пропускания использовались оптические

схемы, приведенные на рис.~\ref{gur:pict2} и~\ref{gur:pict3}. В схеме рис.~\ref{gur:pict2} в качестве

источника света 1 служила двухсотваттная вольфрамовая лампа

накаливания, излучение которой фокусировалось линзой 2 с фокусным

расстоянием 100~мм на входном отверстии болометрического приемника

излучения ПИП-1 5 измерителя мощности ИМО-3 6. 

Непосредственно перед входным отверстием болометрического приемника излучения

ПИП-1 5 располагались светофильтр 3 и измеряемый образец 4. 

Для~перекрытия области спектра от ближней ультрафиолетовой до ближней

инфракрасной было выбрано восемь светофильтров~\cite{gur:bib:4}:

ПС11 (0,320~мкм), ФС6 (0,373~мкм), СС8 (0,410~мкм), ЗС8

(0,529~мкм), ЖС18 (0,600~мкм), ОС14 (0,620~мкм), КС15 (0,700~мкм),

ИКС5 (1,240~мкм).







В этой же схеме измерялись коэффициенты пропускания излучения

полупроводникового лазера с длиной волны 0,630--0,680~мкм.



В схеме рис.~\ref{gur:pict3} в качестве источника света использовался

гелий"=неоновый лазер 1, излучение которого через объектив 2

вводилось в оптоволокно 3, транспортировалось по нему к

измеряемому образцу 4 и регистрировалось с помощью фотодиода ФД-24

5 и вольтметра 6.





\begin{figure}[h!]\centering

\parbox[h]{6.5cm}{\centering \includegraphics[width=6cm]{gur2}

\vspace{0.1mm}\caption{Оптическая схема измерения коэффициентов

пропускания излучения в диапазоне длин волн от ближней

ультрафиолетовой до ближней инфракрасной области спектра: \scriptsize 1 "---

источник света, 2 "--- фокусирующая линза, 3 "--- светофильтр, 4

"--- измеряемый образец, 5 "--- болометрический приемник излучения

ПИП-1, 6 "--- измеритель мощности излучения ИМО-3 \label{gur:pict2}}}\hfill

\parbox[h]{6.5cm}{\centering \includegraphics[width=6cm]{gur3}

\caption{Оптическая схема измерения коэффициентов пропускания

излучения  гелий"=неонового лазера: \scriptsize 1 "--- гелий"=неоновый лазер,

2 "--- объектив, 3 "--- оптоволокно, 4 "--- измеряемый образец, 5

"--- фотодиод ФД-24, 6 "--- вольтметр \label{gur:pict3}}}

\end{figure}



\pagebreak



\Section[n]{Результаты и их обсуждение} Примеры полученных

экспериментальных зависимостей коэффициентов отражения $R$ от длины волны излучения $\lambda $ приведены на рис.~\ref{gur:pict4} и~\ref{gur:pict5}.

Характерной особенностью данных зависимостей является возрастание  коэффициентов отражения в диапазоне длин волн 0,364--0,800~мкм и их последующее уменьшение при $\lambda >0,800$~мкм. 

Выявленные закономерности являются общими, на что, в частности, указывает хорошее совпадение полученных нами результатов для кисти руки человека (рис.~\ref{gur:pict6}) с результатами для кожного покрова человека~\cite{gur:bib:5}, приведенными на рис.~\ref{gur:pict7}.





Столь же общими являются и зависимости коэффициентов пропускания $T$ от длины волны излучения $\lambda $, приведенные на рис.~\ref{gur:pict8} и~\ref{gur:pict9}. 

Эти зависимости указывают на то, что коротковолновое излучение глубже проникает в биологические ткани, но при этом и доля поглощения его оказывается больше. 

Последнее следует из сопоставительного анализа кривых отражения и пропускания. 

Особо обращает на себя внимание сине"=зелёная область спектра, в которой регистрируется возрастание коэффициентов пропускания при одновременном замедлении в росте или даже уменьшении коэффициентов

отражения.



\begin{figure}[p!]\centering

\parbox[h]{6.5cm}{\centering \includegraphics[width=6cm]{gur4}

\vspace{0.1mm}\caption{Зависимости коэффициентов отражения $R$ от

длины волны излучения $\lambda$ для кабачков (К), баклажанов (Б),

помидоров (П), огурцов (О) \label{gur:pict4}}}\hfill

\parbox[h]{6.5cm}{\centering \includegraphics[width=6cm]{gur5}

\caption{Зависимости коэффициентов отражения $R$ от длины волны

излучения $\lambda$ для моркови (М), свеклы (С), картофеля (К),

лука (Л), чеснока (Ч) \label{gur:pict5}}}

%\end{figure}

\smallskip

%\begin{figure}[p!]\centering

\parbox[h]{6.5cm}{\centering \includegraphics[width=5.7cm]{gur6}

\vspace{0.1mm}\caption{Зависимость коэффициента отражения $R$ от

длины волны излучения $\lambda$ для кисти руки человека \label{gur:pict6}}}\hfill

\parbox[h]{6.5cm}{\centering \includegraphics[width=6cm]{gur7}

\caption{Зависимость отражающей способности кожного покрова

человека от длины волны излучения \cite{gur:bib:5}: $A$ "--- слабо

пигментированная кожа, $B$ "--- сильно пигментированная кожа \label{gur:pict7}}}

%\end{figure}

\smallskip

%\begin{figure}[h!]\centering

\parbox[h]{6.5cm}{\centering \includegraphics[width=6cm]{gur8}

\vspace{0.1mm}\caption{Зависимости коэффициентов пропускания $T$

от длины волны излучения $\lambda$ для свеклы толщиной 0,5~мм (1);

1,0~мм (2) и 1,5~мм (3) \label{gur:pict8}}}\hfill

\parbox[h]{6.5cm}{\centering \includegraphics[width=6cm]{gur9}

\caption{Зависимости коэффициентов пропускания $T$ от длины волны

излучения $\lambda$ для моркови толщиной 0,5~мм (1); 1,0~мм (2) и

1,5~мм (3) \label{gur:pict9}}}

\end{figure}





Для одних и тех же биологических тканей коэффициенты пропускания

возрастают с увеличением мощности излучения (рис.~\ref{gur:pict10}) и

уменьшением оптической плотности (разрыхлением) биологических

тканей (рис.~\ref{gur:pict11}).



Анализ особенностей поведения кривых отражения и пропускания

позволяет предположить, что в их основе лежат глубокие физические

причины, обусловленные влиянием излучения Солнца на биосферу

Земли. Действительно, к настоящему времени собрано достаточное

количество фактов [\citen{gur:bib:6}--\citen{gur:bib:9}],

свидетельствующих о том, что любой живой организм, как и вся

биосфера Земли в целом, по"=своему служат избирательными

резонаторами и трансформаторами солнечной активности. Жизнь на

Земле возникла и развивалась благодаря солнечной энергии в течение

нескольких миллиардов лет. Поэтому, как справедливо впервые

отметил профессор А.~Л.~Чижевский, каждая живая клетка по

существу представляет собой своеобразный трансформатор,

преобразующий энергию Солнца в тепловую, механическую,

электрическую, химическую и другие энергии. 

И~именно спектр излучения Солнца, под действием которого сформировались устойчивые жизненные циклы клеток, должен проявляться в избирательном характере процессов взаимодействия световых потоков с~биологическими тканями.



\begin{figure}[h!]\centering

\parbox[h]{6.5cm}{\centering \includegraphics[width=6cm]{gur10}

\vspace{0.1mm}\caption{Зависимости коэффициентов пропускания $T$

от толщины $d$ картофеля для различных источнико света в красной

области спектра: \scriptsize 1 "--- вольфрамовая лампа накаливания, 2 "---

полупроводниковый лазер, 3 "--- гелий"=неоновый лазер \label{gur:pict10}}}\hfill

\parbox[h]{6.5cm}{\centering \includegraphics[width=6cm]{gur11}

\caption{Зависимости коэффициентов пропускания $T$ от толщины $d$

картофеля и моркови для излучения гелий"=неонового лазера: \scriptsize  1 "---

сырой картофель, 2 "--- печёный картофель, 3 "--- сырая морковь, 4

"--- подсушенная морковь, 5 "--- сухая морковь \label{gur:pict11}}}

\end{figure}





Спектр излучения Солнца \cite{gur:bib:10} приведен на рис.~\ref{gur:pict12}.

Видно, что его максимум приходится на сине"=зелёную область. 

Это в

соответствии с законом смещения Вина \cite{gur:bib:11}

\begin{equation}

\label{gur:eq1} T_\text{ц} \lambda_{\max} =b,

\end{equation}

где $T_\text{ц}$ "--- цветовая температура, $\lambda_{\max}$ "---

длина волны, при которой наблюдается максимум спектра излучения,

$b \approx {hc}/({4,965k})=2,898\cdot 10^{-3}$~м$\cdot $К "---

постоянная Вина; $h=6,626\cdot 10^{-34}$~Дж$\cdot $с "---

постоянная Планка; $c=2,998\cdot 10^8$~м/с "--- скорость света в

вакууме; $k=1,380\cdot 10^{-23}$~Дж/К "--- постоянная Больцмана,

позволяет оценит цветовую температуру на поверхности Солнца $\sim

6\cdot 10^3$~К ($\lambda _{\max} \approx 0,480$~мкм).





\begin{figure}[t!] \centering

\includegraphics{gur12}

\caption{Спектр излучения солнца~\cite{gur:bib:10} \label{gur:pict12}}

\end{figure}



Отметим, что в отличие от спектра излучения Солнца максимум спектра излучения вольфрамовой лампы накаливания, использованной нами в схеме рис.~\ref{gur:pict2}, приходится на инфракрасную область (рис.~\ref{gur:pict13}).

На рис.~\ref{gur:pict13} наилучшее согласие экспериментальных данных с расчётным~спектром излучения по формуле Планка~\cite{gur:bib:11}, нормированной на единицу

\begin{equation}

\label{gur:eq2} I=\frac{\lambda _{\max}^5 \left[\exp \left(

\frac{hc}{bk}\right)-1\right]}{\lambda^5\left[ \exp \left(

\frac{hc}{bk}\frac{\lambda _{\max} }{\lambda } \right)-1 \right]},

\end{equation}

достигается при $T_\text{ц} \approx 3,6\cdot 10^3$~К ($\lambda_{\max} \approx 0,800$~мкм).



В годы повышенной солнечной активности, сопровождаемой образованием на Солнце локальных участков с~температурой, значительно превышающей $6 \cdot 10^3$~К, максимум спектра излучения Солнца смещается в~коротковолновую область (рис.~\ref{gur:pict14}).

Мощные потоки коротковолнового (ультрафиолетового и~рентгеновского) излучения совместно с~потоком заряженных и~электрически нейтральных частиц на время выводят биосферу Земли из состояния равновесия. Эволюционный процесс развития под благотворным влиянием солнечной энергии прерывается всплеском энергетической активности, что, как установил профессор А.~Л.~Чижевский, может служить одной из причин возникновения стихийных бедствий: наводнений, засух, землетрясений, извержений вулканов, массовых нашествий вредных насекомых, эпидемий, эпизотий (эпидемических заболеваний животных), эпифитий (эпидемических заболеваний растений), обострений сердечно"=сосудистых и~онкологических заболеваний и~даже социальных конфликтов в~обществе. 

Периодическое (с~периодом $\sim  11,1$~года) возрастание солнечной активности почти повсеместно приводит к~особому изменению состава крови, к уменьшению в ней белых кровяных клеток "--- лейкоцитов и одновременному увеличению содержания в~крови лимфоцитов, т.~е. наблюдается такое изменение состава крови, как после радиоактивного облучения. 

Причём кровь сильнее реагирует на возрастание солнечной активности по мере приближения к полюсам Земли, где атмосфера становится особенно тонкой и проницаемой (возникают даже озоновые дыры) для солнечной радиации.





\begin{figure}[t!]\centering

\parbox[h]{6.5cm}{\centering \includegraphics[width=6cm]{gur13}

\vspace{0.1mm}\caption{Экспериментальный и рассчитанные по

нормированной на единицу формуле Планка спектры

излучения вольфрамовой лампы накаливания: \scriptsize 1 "--- $T_\text{ц}

\approx 3,6\cdot 10^3$~K, $\lambda_{\max} \approx 0,800$~мкм; 2

"--- $T_\text{ц} \approx 3,2\cdot 10^3$~K, $\lambda _{\max }

\approx 0,900$~мкм; 3 "--- $T_{\text{ц}} \approx 2,9\cdot 10^3$~K,

$\lambda_{\max}\approx 1,000$~мкм \label{gur:pict13}}}\hfill

\parbox[h]{6.5cm}{\centering \includegraphics[width=6cm]{gur14}

\caption{Рассчитанные по приведённым в тексте формулам зависимости от температуры Солнца $T_c$ длины волны

$\lambda_{\max}$, соответствующей максимуму спектра излучения

Солнца (1), средней температуры Земли $T_\text{з}$

(2) и среднего радиуса земной орбиты $R$

(3)} \label{gur:pict14}}

\end{figure}





Солнечное излучение лежит в основе реакции <<Ф>> в крови, суть которой состоит в выпадении белковых хлопьев в сыворотке крови под воздействием определенных реактивов. 

Установлена прямая зависимость реакции <<Ф>> от суточного хода Солнца, т.~е. от интенсивности его излучения. 

Реакция <<Ф>> уменьшается во время солнечных затмений как полных, так и частичных. 

При этом кровь почти мгновенно реагирует на изменение интенсивности солнечного излучения при затмении.



Таким образом, из анализа представленных результатов следует, что влияние световых потоков на биологические ткани в период их устойчивого развития, т.~е. при нахождении биологической системы в~равновесном состоянии, будет благотворным лишь при интенсивностях, сопоставимых с интенсивностью излучения спокойного Солнца в~соответствующем интервале длин волн. 

Если же биологическая система по каким"=либо причинам оказывается выведенной из состояния равновесия, то для возврата её в~исходное равновесное состояние могут быть использованы световые потоки большей интенсивности и~преимущественно сине"=зеленой области спектра. 

При этом чем больше интенсивность используемого излучения, тем меньше должно быть время его благотворного воздействия для стимулирования механизма самовозврата биологической системы в~состояние равновесия.

Превышение данного времени воздействия чревато раскачиванием биологической системы, что может проявиться в ослаблении защитных функций организма.



По этим причинам весьма проблематично использование лазерного излучения для повышения урожайности сельскохозяйственных культур только на этапе облучения их семян, тогда как в~остальной период своего развития растения будут находиться под воздействием лишь солнечного света. 

Более того, при таком подходе лазерное излучение должно оказывать гнетущее воздействие на развитие растений, что, как правило, и подтверждается практикой возврата к исходной или даже снижения урожайности сельскохозяйственных культур уже во втором их поколении~\cite{gur:bib:3}.



Можно также утверждать, что использование в лазерной терапии излучения гелий"=неонового лазера не является оптимальным. 

Для целей терапии наибольшей эффективностью должны обладать лазеры, длины волн излучения которых приходятся на сине"=зелёную область спектра. 

Как видно из данных таблицы, к числу таких лазеров, в~частности, относятся компактные полупроводниковые лазеры ZnSe, CdS, ZnTe, из которых лазер CdS может работать при комнатной температуре, а~также аргоновый лазер.



\begin{table}[b!]



\vspace{-3mm}



\centering \small \caption*{\bf Лазеры с длинами волн излучения

от~ближней ультрафиолетовой до~ближней инфракрасной области

спектра~\cite{gur:bib:5,gur:bib:12}} \vspace{-3mm}

\begin{tabular}{c|c|c} \hline

Тип лазера& Тип рабочей среды & Длина волны излучения, мкм

\\\hline

Газовые  & Аргоновый лазер & 0,488; 0,515 \\ \cline{2-3}

лазеры   & Гелий"=неоновый лазер& 0,633 \\ \hline

Твердотельные &  Рубиновый лазер & 0,694 \\ \cline{2-3}

лазеры        &  Неодимовый лазер& 1,060 \\ \hline

  & ZnS-лазер & 0,330 \\ \cline{2-3}

                         & ZnO-лазер & 0,370 \\ \cline{2-3}

                         & ZnSe-лазер& 0,460 \\ \cline{2-3}

                         & CdS-лазер & 0,490 \\ \cline{2-3}

Полупроводниковые                         & ZnTe-лазер& 0,530 \\ \cline{2-3}

лазеры                         & CaSe-лазер& 0,590 \\ \cline{2-3}

                         & CdSe-лазер& 0,675 \\ \cline{2-3}

                         & CdTe-лазер& 0,785 \\ \cline{2-3}

                         & GaAs-лазер& 0,840--0,950 \\ \cline{2-3}

                         & InP-лазер & 0,910 \\ \cline{2-3}

                         & GaSb-лазер& 1,550 \\ \cline{2-3}

                         & InAs-лазер& 3,100 \\ \hline

\end{tabular}

\end{table}



С точки зрения перспектив развития биосферы Земли под воздействием излучения Солнца, которое неизбежно рано или поздно угаснет, но перед своим угасанием должно будет пройти этап вспышки, в~соответствии с~данными рис.~\ref{gur:pict14} в~медико"=биологических исследованиях важное место должно быть отведено изучению длительного воздействия коротковолнового излучения на формирование новых устойчивых состояний биологических систем, их новых жизнестойких форм. 

При этом следует также иметь в виду, что с повышением температуры Солнца средняя температура Земли будет повышаться как

\begin{equation}

\label{gur:eq3}

T_\text{З} (\text{\textcelsius})=T_\text{С} ({\rm K})\sqrt {\frac{R_\text{С}}{2R}} -273,

\end{equation}

где $R_\text{С} =6,960\cdot 10^8$~м "--- радиус Солнца, $R=1,496\cdot 10^{11}$~м "--- средний радиус земной орбиты, $T_\text{С}$ "--- температура Солнца в градусах Кельвина, $T_\text{З}$ "--- средняя температура Земли в градусах Цельсия. 

Чтобы средняя температура Земли с повышением температуры Солнца оставалась такой же, какой она является сейчас, $\sim 290$~К или $\sim 17~$\textcelsius, необходимо, чтобы средний радиус земной орбиты возрастал как

\begin{equation}

\label{gur:eq4}

R=R_\text{С} \frac{T_\text{С} ^2(K)}{2T_\text{З}^2(K)}.

\end{equation}

Осуществление этого на практике представляется столь же маловероятным, как и~организация защиты биосферы Земли от воздействия коротковолнового излучения.

А это означает, что биосфера Земли под воздействием коротковолнового излучения и~температуры должна будет неизбежно эволюционировать в~новое устойчивое состояние.



\smallskip



\Section[n]{Выводы} Таким образом, экспериментально изучены особенности

взаимодействия световых потоков с биологическими тканями. На основе

полученных результатов и современных представлений о влиянии Солнца на

биосферу Земли дано физическое обоснование оптимального выбора лазерных

источников и условий их практического применения для решения фундаментальных

и прикладных задач медико"=биологических исследований.



\begin{thebibliography}{99}



\RBibitem{gur:bib:1}

\book Лазеры в клинической медицине

\ed проф. С.~Д.~Плетнёв

\publaddr М.

\publ Медицина

\yr 1981

\totalpages 400

\misctransl

\ed Prof. S.~D.~Pletnev

\book Lasers in Clinical Medicine

\yr 1981

\totalpages 400

\publ Medicine

\publaddr Moscow



\RBibitem{gur:bib:2}

\by H.-P.~Berlien, G.~M\"uller

\publ Ecomed

\publaddr Landsberg

\yr 1996 

\book Angewandte Lasermedizin

\bookinfo Lehr- und Handbuch f\"ur Praxis und Klinik

\rtransl русск. пер.:~

\book Прикладная лазерная медицина

\bookinfo Учебное и справочное пособие

\by  Х.-П.~Берлиeн, Г.~Мюллер 

\publaddr  М.

\publ Интерэксперт

\yr 1997





\RBibitem{gur:bib:3}

\by  В.~Г.~Володин, В.~А.~Мостовников, Б.~И.~Авраменко, И.~В.~Хохлов

\book Лазеры и наследственность растений

\publaddr Минск

\publ Наука и техника

\yr 1984

\totalpages 175

\misctransl

\by  V.~G.~Volodin, V.~A.~Mostovnikov, B.~I.~Avramenko, I.~V.~Khokhlov

\book Lasers and Plant Genetics 

\publ Nauka i Tekhnika

\publaddr Minsk 

\yr 1984

\totalpages 175





\RBibitem{gur:bib:4}

\by А.~Н.~Зайдель, Г.~В.~Островская, Ю.~И.~Островский

\book Техника и практика спектроскопии

\serial Физика и техника спектрального анализа

\publaddr М.

\publ Наука

\yr 1976

\totalpages 392

\misctransl

\by A.~N.~Zaidel, G.~V.~Ostrovskaya, Yu.~I.~Ostrovsky

\book Techniques and Practices of Spectroscopy

\publ Nauka

\yr 1976 

\publaddr Moscow

\totalpages 392





\RBibitem{gur:bib:5}

\by Реди~Дж.

\book Промышленное применение лазеров

\publaddr М.

\publ Мир

\yr 1981

\totalpages 640



\RBibitem{gur:bib:6}

\by А.~Л.~Чижевский

\book Солнце и мы

\publaddr М.

\publ Знание

\yr 1963

\totalpages 48

\misctransl

\by A.~L.~Chizhevskiy

\book The Sun and us

\publaddr Moscow

\publ Znanie

\yr 1963

\totalpages 48



\RBibitem{gur:bib:7}

\by А.~Л.~Чижевский, Ю.~Г.~Шишина

\book В ритме Солнца

\publaddr М.

\publ Наука

\yr 1969

\totalpages 112

\misctransl

\book In the Rhythm of the Sun

\by A.~L.~Chizhevskiy, Yu.~G.~Shishina 

\publaddr Moscow \publ Nauka \yr 1969 \totalpages 112





\RBibitem{gur:bib:8}

\by А.~Л.~Чижевский

\book Земное эхо солнечных бурь

\publaddr М.

\publ Мысль

\yr 1976

\totalpages 366

\misctransl

\by \by  A.~L.~Chizhevsky

\book  The Terrestrial Echo of Solar Storms

\publaddr Moscow

\publ Mysl

\yr 1976

\totalpages 366



\RBibitem{gur:bib:9}

\by А.~Л.~Чижевский

\book Космический пульс жизни. Земля в объятиях Солнца. Гелиотараксия

\publaddr М.

\publ Мысль

\yr 1995

\totalpages 766

\misctransl

\book Space Pulse of Life. The Earth in the Embrace of the Sun. Geliotaraksiya

\publaddr Moscow

\publ Mysl

\by  A.~L.~Chizhevsky

\yr 1995

\totalpages 766



\RBibitem{gur:bib:10}

\book Физические величины

\bookinfo Справочник

\eds И.~С.~Григорьев, Е.~З.~Мейлихов

\publaddr М.

\publ Энерго\-атом\-издат

\yr 1991

\totalpages 1232

\misctransl

\book Physical Quantities

\bookinfo Handbook

\eds I.~S.~Grigoriev, E.~Z.~Meylikhov

\publaddr Moscow

\publ Energoatomizdat

\yr 1991

\totalpages 1232





\RBibitem{gur:bib:11}

\by Г.~С.~Ландсберг

\book Оптика

\publaddr М.

\publ Наука

\yr 1976

\totalpages 928

\misctransl

\by G.~S.~Landsberg  

\yr 1976 

\book Optics 

\publaddr Moscow

\publ Nauka

\totalpages 928



\RBibitem{gur:bib:12}

\book Справочник по лазерам

\bookinfo в~2-х томах

\bookvol 1

\ed А.\,М.~Прохоров

\publaddr М.

\publ Советское радио

\yr 1978

\totalpages 504

\misctransl

\book Handbook of Lasers

\ed A.~M.~Prokhorov

\publaddr Moscow

\publ Sovetskoe Radio

\yr 1978

\totalpages 504

\bookvol 1



\end{thebibliography}







\noindent

\hskip6.5cm \thedate \par\noindent

\hskip6.5cm\therevdate







%\chead[\fancyplain{}{\scriptsize \sl D.~M.~G\,u\,r\,e\,e\,v}]{\fancyplain{}{\scriptsize \sl  Medical and biological aspects of laser influence}}



\makeengtitlem



\newpage



%\end{document}

